\noindent\textbf{Background.} Algebraic connectivity $\lam_2$ and its closed-form
per-edge sensitivity govern stability in coupled excitable networks, motivating
a biomarker of atrial reentry, the Spectral Fragility Index (\SFI). Whether it
adds information beyond established fibrosis and connectivity features was
unknown.

\medskip
\noindent\textbf{Methods.} Under a protocol pre-registered before any label was
inspected, \SFI\ was added to six competitor features on $182$ patient-derived
left-atrial meshes from two public cohorts (Roney, $n=100$; University of
Washington, $n=82$) and up to $100{,}000$ synthetic excitable networks. The
in-silico outcome was a simulator's reentry-inducibility verdict; a clinical arm
used two-year post-ablation arrhythmia recurrence in the $82$ Washington
patients. Models were fitted under patient- and shape-family-held-out
cross-validation and compared by paired DeLong test. The primary endpoint
required grouped $\dAUC\geq0.05$ with $p<0.05$.

\medskip
\noindent\textbf{Results.} The primary endpoint was not met. On the combined
cohorts ($42$ inducible of $182$), gradient boosting gave $\dAUC=+0.012$
(bootstrap CI $-0.017$ to $+0.044$; $p=0.449$), an interval that excludes the
$0.05$ threshold; synthetic-network $\dAUC$ converged to $+0.003$. The clinical
gate fired ($\dAUC=+0.197$; $p=0.0061$), but the pre-specified plan identified
an artefact of an anti-predictive baseline (AUC $0.250$) that nine columns of
Gaussian noise reproduced (empirical $p=0.333$); that endpoint was not met. The
estimator was ill-conditioned on real tissue (median $\rhoval\approx2422$;
validity limit $\rhoval^{\ast}\approx3$).

\medskip
\noindent\textbf{Conclusions.} \SFI\ added no practically meaningful predictive
value for atrial reentry beyond existing features, and its one nominally
positive patient-outcome result was a false positive its own pre-registered
guards caught. These findings do not support clinical use.
